% Part:sets-functions-relations
% Chapter: size-of-sets
% Section: introduction

\documentclass[../../../include/open-logic-section]{subfiles}

\begin{document}

\olfileid{sfr}{siz}{int}

\olsection{Inleiding}

Toen Georg Cantor in de jaren 1870 de verzamelingenleer opbouwde, wilde hij het
idee van een oneindige verzameling aanvaardbaar maken. De
middeleeuwers zouden dit een actuele oneindigheid noemen: een oneindigheid die
als voltooid wordt opgevat. Daarvoor keek hij vooral naar de \emph{grootte} van
verzamelingen. Als $a$, $b$ en $c$
drie verschillende dingen zijn, voelt de verzameling $\{a, b, c\}$
\emph{groter} aan dan $\{a, b\}$. Maar hoe werkt dat bij oneindige
verzamelingen? Zijn die allemaal even groot? Nee, zo blijkt.

Het eerste idee dat we nodig hebben is een opsomming: we zetten de
!!{element}s één voor één in een lijst. Bij elke eindige verzameling kan dat.
Bij sommige oneindige verzamelingen kan het ook, zolang de lijst zelf maar
oneindig mag zijn. Zulke verzamelingen noemen we !!{enumerable}. Cantor kwam
tot een verrassende uitkomst die aan het einde van dit hoofdstuk helemaal
duidelijk zal zijn: sommige oneindige verzamelingen zijn niet !!{enumerable}.
\end{document}
